\documentclass[11pt]{article}
\usepackage[a4paper,margin=26mm]{geometry}
\usepackage{amsmath,amssymb,amsthm,mathtools}
\usepackage[T1]{fontenc}
\usepackage{lmodern,microtype}
\usepackage[colorlinks=true,linkcolor=blue,urlcolor=blue,citecolor=blue]{hyperref}
\newtheorem{theorem}{Theorem}
\newtheorem{lemma}[theorem]{Lemma}
\newtheorem{proposition}[theorem]{Proposition}
\theoremstyle{remark}\newtheorem*{remark}{Remark}
\newcommand{\E}{\mathbb E}
\newcommand{\Prob}{\mathbb P}
\newcommand{\eps}{\varepsilon}
\newcommand{\NN}{\mathcal N}
\newcommand{\Own}{\mathsf O}
\newcommand{\Lost}{\mathsf L}
\DeclareMathOperator{\dist}{dist}
\title{A shorter route to the Erd\H{o}s--Gallai conjecture}
\author{Working reconstruction of the proof claim by Ryan Coffey}
\date{1 October 2026}
\begin{document}
\maketitle
\begin{abstract}
We give a substantially shorter reconstruction of the claimed proof: one
kind of expander piece, one retirement argument, elementary grouping of
two-edge paths, and a small quotient. A bounded-degree auxiliary graph
replaces the long case distinction in the random linking argument.
The main external combinatorial inputs are Haxell's matching theorem
and the Lov\'asz path decomposition corollary.
This is a new mathematical draft, checked by static reasoning but not
formally verified. It is not a transcription of the submitted Lean proof.
\end{abstract}

\section{The idea}
Write $f(G)$ for the least number of simple cycles and single edges whose
edge sets partition $E(G)$. All graphs are finite and simple, all logarithms
are base two, and all implicit constants are absolute.

\begin{theorem}\label{thm:main}
There is an absolute constant $C$ such that $f(G)\le C|V(G)|$.
\end{theorem}

The useful target is a reduction, rather than a decomposition in one step:
\begin{equation}\label{eq:reduction}
 f(G)\le C_0 n+2\sum_i f(Q_i),\qquad
 \sum_i|V(Q_i)|\le n/4.
\end{equation}
The $Q_i$ are simple graphs. Induction then proves the theorem with
$C=2C_0$. The factor two only pays for quotient edges representing paths
of two original edges.

We repeatedly extract robust expanders, with the average degree of the
remaining graph falling from $d$ to a fixed power of $\log d$. A vertex
is new in at most two rounds. Vertices shared between pieces have small
total incidence. These are the vertices we pay to remove locally.
Every other vertex can use an expander from at least two rounds earlier
as a reservoir of connecting edges.

Most remaining edges close into cycles by a particularly simple rule:
pair edges into two-edge paths, group vertex-disjoint paths greedily,
and connect each group cyclically. The unpaired edges have at most one
edge at each pair consisting of a shared vertex and an earlier piece.
This property lets us move their endpoints into tiny random pools and
obtain the graphs $Q_i$ in \eqref{eq:reduction}.

The order matters. We build the hierarchy forwards, but use its pieces
backwards. A reservoir is lent before its owner is decomposed. Only edges
actually used by the borrower are removed from the owner's graph.

\section{Expansion and sparse linking}
An $N$-vertex graph is \emph{$(\eps,s)$-robust} if
\[
 |\NN_{G-F}(U)|\ge \frac{\eps|U|}{\log^2 N}
 \quad\text{whenever }1\le |U|\le 2N/3,\quad |F|\le s|U|.
\]
Here $\NN$ is the external neighbourhood. Graphs of order at most one
are treated as terminal pieces. In the hierarchy $\eps=1/32$.

We use the following corollary of Lov\'asz's theorem, as stated in
\cite[Corollary 22]{BM}: every graph has a decomposition into simple
paths with each vertex an endpoint at most twice.
We also use Haxell's theorem \cite{Haxell} in this form: families
$\mathcal H_i$ of sets of size at most $h$ have pairwise disjoint
representatives if, for every nonempty index set $I$ and every set $F$
with $|F|\le(2h-1)(|I|-1)$, some member of some $\mathcal H_i$, $i\in I$,
avoids $F$. Chernoff and martingale Bernstein inequalities are used below.

\begin{lemma}[Edge splitting]\label{lem:split}
Suppose edges of an $(\eps,s)$-robust $N$-vertex graph independently
receive labels, each specified label having probability $p$.
If $ps\ge C\log N$, its label graph is $(\eps,ps/2)$-robust with
failure probability at most $N^{-6}$. The constant $C$ can be increased
to give any fixed polynomial failure bound.
\end{lemma}
\begin{proof}
For a set $U$ of size $u\le2N/3$ and an external set $T$ of size
$|T|<\eps u/\log^2N$, robustness forces more than $su$ edges from $U$
to $V\setminus(U\cup T)$. Chernoff bounds the probability that at most
$psu/2$ have the label by $\exp(-psu/8)$. There are at most $N^{3u}$
choices of $(U,T)$, up to an inessential constant. Sum over $u$.
Failure of the desired robustness supplies exactly such a pair.
\end{proof}

The following is the main analytic simplification. Its proof is included
so that no source-specific linking theorem is left as a black box.

\begin{lemma}[Sparse linking]\label{lem:link}
Fix $\eps_0>0$. There are constants $C,N_0$ depending only on $\eps_0$
with the following property. Let $G$ be $(\eps,s)$-robust on $N\ge N_0$
vertices, where $\eps_0\le\eps\le1$. Put $L=\log N$. If
\[
 0<\rho\le\tfrac12,\qquad t\ge1,\qquad
 s\ge CtL^{18}\rho^{-3},
\]
then a $\rho$-random vertex set $V$ has, with probability at least
$1-N^{-3}$, this property: every multiset of pairs of distinct vertices,
with each vertex occurring at most $t$ times, can be joined by
edge-disjoint paths of length $O(L^4)$ whose interiors lie in $V$.
The pairs may be chosen after $V$.
\end{lemma}
\begin{proof}
Put $\alpha=\eps/L^2$. A ball $\mathcal B_G^a(U;V)$ through $V$
counts endpoints in $V$ reachable from $U$ by paths of length at most
$a$ whose interiors lie in $V$; starting vertices are unrestricted,
and zero-length paths are allowed.

\emph{A short certificate for random growth.}
Let $1\le|W|\le2N/3$, $|F|\le s|W|/4$, and $1\le d\le s$ be an
integer. Take a maximal bipartite subgraph $H$ of $G-F$ across
$(W,V(G)\setminus W)$ with degree at most $d$ on both sides. Then
\begin{equation}\label{eq:certificate}
 e(H)\ge\alpha d|W|/4.
\end{equation}
Indeed, otherwise the sets $S\subseteq W$ and $X\subseteq V(G)\setminus W$
of saturated vertices each have size less than $\alpha|W|/4$.
For $U=W\setminus S$, maximality gives
$\NN_{G-F-H}(U)\subseteq S\cup X$. But $|U|\ge3|W|/4$ and
$|F\cup E(H)|\le s|W|/2\le s|U|$, contradicting expansion.

Sample vertices of $W$ independently with probability $p$, and take
$d=\lceil1/p\rceil\le s$. Let $Z$ count their $H$-neighbours outside $W$.
Since $pr\le2$ for every right degree $r$, we have
\[
 \mu:=\E Z\ge pe(H)/3\ge\alpha|W|/12.
\]
Expose the independent vertex bits. The resulting Doob martingale has
increments bounded by $d$ and conditional variance sum at most
\[
 \beta=\sum_{w\in W}p(1-p)\deg_H(w)^2
 \le pd e(H)\le3d\mu.
\]
Bernstein, with deviation $\mu/2$, therefore gives
\begin{equation}\label{eq:growth}
 \Prob(Z<\alpha|W|/24)\le\exp(-\alpha p|W|/624).
\end{equation}
No independence among the right-side coverage events is needed.

\emph{Sprinkling.}
For this paragraph suppose the robustness parameter is at least
$\theta+1$, where
\[
 \ell=\lceil100L/\alpha\rceil,\qquad p=\rho/(10\ell),\qquad
 \theta=C_1L/(\alpha p),\qquad \mu_0=\alpha/(6\theta).
\]
Thus $\ell=O(L^3)$, $\theta=O(L^6/\rho)$ and
$\mu_0^{-1}=O(L^8/\rho)$.
Call $A$ well-expanding if $|\NN_G(A)|\ge\theta|A|$.
For fixed such $A$, $a=|A|>0$, and $|F|\le a$, expose $\ell-1$
independent $p$-random batches followed by one independent $q$-random
batch, chosen so their union has density exactly $\rho$.
The earlier union has density at most $\rho/10$, so $q\ge0.9\rho$.

Let $B_i$ be the vertices reachable from $A$ by paths of length at most
$i$, with interiors in the first $i-1$ batches and unrestricted endpoints;
include zero-length paths. Initially $|B_1|\ge\theta a$.
Conditional on the past, \eqref{eq:growth} shows that, while
$|B_i|\le2N/3$, the next batch increases it by at least
$\alpha|B_i|/24$, except with probability
$\exp(-\alpha p\theta a/624)$.
All hypotheses hold since $|F|\le a$ and the robustness is at least
$\theta+1$. Choosing $C_1$ large makes the total failure at most
$N^{-11a}$. The $\ell-1$ growth steps suffice to pass $2N/3$.
The last batch then contributes more than $0.56\rho N$ reachable
vertices, while the full union has fewer than $1.1\rho N$, except with
probability $\exp(-c\rho N)$. Since $\theta a\le N$ and
$\rho\theta\gg C_1L$, this is at most $N^{-11a}$.
Consequently the ball of radius $\ell$ from $A$, restricted to endpoints
in $V$, has size greater than $|V|/2$, with failure at most $N^{-10a}$.

Every $U$ with $1\le|U|\le2N/3$ contains a well-expanding
$A$ of size at least $\alpha|U|/(3\theta)$.
Take an inclusion-maximal well-expanding $A\subseteq U$. If $A=U$, this is immediate.
Otherwise every $u\in U\setminus A$ has fewer than $\theta+1$
neighbours outside $A\cup\NN(A)$. Delete these edges. Expansion gives
$\alpha|U|\le|\NN(A)|$, so $A\ne\varnothing$; maximality also gives
$|\NN(A)|<\theta(|A|+1)+1\le3\theta|A|$.
This proves the assertion, including $A\ne\varnothing$.

A union bound over the $O(N^{3a})$ choices of $(A,F)$ with
$|A|=a$, $|F|\le a$, now gives the following event with failure at most
$N^{-6}$:
\begin{equation}\label{eq:balls}
 |\mathcal B_{G-F}^{\ell}(U;V)|>|V|/2
 \quad\text{for all }U\ne\varnothing,\ |F|\le\mu_0|U|.
\end{equation}
For $|U|>2N/3$, use a subset of size between $N/2$ and $2N/3$;
the factor six in $\mu_0$ covers this case.

\emph{From balls to paths.}
Here is a useful elementary observation from the argument of \cite{BM}.
Suppose every $u$-vertex set reaches more than half of $V$ in radius
$\ell$ and $|V|\ge4u-2$. At most $u-1$ vertices can have a ball of
radius $2\ell L$ of size at most $|V|/2$. Otherwise start with a
$u$-set of such vertices. Split it into two almost equal sets; one
half's current ball contains at least $u$ vertices, since the union
contains more than $|V|/2\ge2u-1$. Expanding those $u$ vertices by
another $\ell$ steps restores a ball larger than half of $V$.
Repeat while the starting set has size at least two, until a singleton
remains, in at most $\lceil\log u\rceil$ steps. Its final radius is at most
$(1+\lceil\log u\rceil)\ell\le2\ell L$ for $N\ge4$.
This contradicts the choice of the singleton; $u=1$ is immediate. Concatenation is
legitimate because the intermediate ball vertices belong to $V$.
It follows that among $2u-1$ disjoint pairs some pair has a joining
path through $V$ of length at most $h=\lceil4\ell L\rceil$.

Suppose \eqref{eq:balls} holds with deletion allowance
$\bar s|U|$, where $\bar s=40ht/\rho$, and $|V|\ge\rho N/2$.
For a nonempty submultiset $I$ of demands, choose a maximal collection
of vertex-disjoint pairs, of size $q$; then $|I|\le2tq$.
Put $M=\lceil4/\rho\rceil$ and $u=\lceil q/(2M)\rceil$.
We have $2u-1\le q$ and, since $q\le N/2$,
$4u-2\le N/M+2\le\rho N/2\le|V|$ when $\rho N\ge8$.
Any forbidden edge set satisfying Haxell's threshold obeys
\[
 |F|\le(2h-1)(|I|-1)<2h|I|\le4htq\le40htu/\rho.
\]
The observation supplies a path avoiding $F$ for one of these pairs.
Apply Haxell to the edge sets of the length-at-most-$h$ joining paths.
This gives all the desired edge-disjoint paths, including repeated demands.

\emph{Amplification.}
Set $K=\lceil\bar s/\mu_0\rceil=O(tL^{12}\rho^{-2})$ and split the
edges uniformly into $K$ colours. Lemma~\ref{lem:split}, with a stronger
fixed polynomial tail if needed, gives robustness $s/(2K)\ge\theta+1$
in every colour. Apply \eqref{eq:balls} in each colour, using the same
independent set $V$. Given $|F|\le\bar s|U|$, some colour contains at
most $\mu_0|U|$ edges of $F$, so its large ball proves the amplified
property for $G-F$. It suffices that
$s\ge2K(\theta+1)=O(tL^{18}\rho^{-3})$.
Feasibility gives $N>s$, hence $K<N$ and $\rho N\gg\log N$.
The splitting failure, $KN^{-6}$, and the Chernoff failure for
$|V|\ge\rho N/2$ total at most $N^{-3}$ for large $N$.
The auxiliary colouring merely witnesses a property of $(G,V)$, so it
can be discarded. This proves the uniform assertion of the lemma.
\end{proof}

\section{A hierarchy with only one kind of piece}
Fix a sufficiently large absolute cutoff $D$. Keep the original $n$
vertices throughout. In a round with residual average degree $d\ge D$, put
\begin{equation}\label{eq:scales}
 \lambda=\log d,\quad M=\lceil d^2\rceil,\quad
 s=\lceil\lambda^{100}\rceil,\quad P=\lceil\lambda^{103}\rceil,
 \quad\tau=\lceil2^{14}\lambda^{102}\rceil.
\end{equation}
The exponents provide slack; their numerical values are not significant.

\emph{The splitting rule.}
If a graph $H$ of order $m$ is not $(\eps,s)$-robust, take a witness
$U,F$, put $N=\NN_{H-F}(U)$ and
$F_0=E_H(U,V(H)\setminus(U\cup N))$.
Move the vertices of $U$ with $F_0$-degree at least $\tau$ to $N$;
call the remaining set $A$. Then also move to $N$ every vertex outside
$U\cup N$ with at least $\tau$ neighbours in $A$.
Call the resulting separator $B$, and delete the edges
$F''=E_H(A,V(H)\setminus(A\cup B))$.
The children are
\[
 H[A\cup B],\qquad H[V(H)\setminus A]-E_H(B).
\]
Their edge sets and $F''$ partition $E(H)$; their vertex intersection is $B$.
If $\tau\ge256s\log^2m_0$, where $m_0$ is the root order, then
\begin{equation}\label{eq:splitshape}
 |B|\le\frac{2\eps|A|}{\log^2m},\quad
 0<|A|,\quad |A\cup B|<3m/4,\quad |F''|\le2s|A|.
\end{equation}
Indeed the edges charged to the two kinds of moved vertices are disjoint
subsets of $F_0$. Their total number of vertices is at most
$s|U|/\tau\le\eps|U|/(8\log^2m)$; combine this with the witness bound
on $N$. Both children are smaller, so recursion terminates. For $s=0$
there are no deletions or moved vertices and the same shape bounds hold.

\emph{Overlap and thin cuts.}
In any such splitting tree let $S$ be the sum of leaf orders and
$\Delta$ the sum of separator sizes. Exact counting gives $S=n+\Delta$.
At a split of order $m$, charge $2\eps/\log^2m$ for each $u\in A$
to a leaf copy of $u$ below the first child. Along the nodes charging
one copy, orders grow by a factor at least $4/3$. Hence separators at
nodes of order at least $T\ge2$ have total size at most
\[
 2\eps S\sum_{j\ge0}(\log T+j\log(4/3))^{-2}
 <\frac{7\eps S}{\log T}.
\]
Taking $T=2$ and $\eps=1/32$ gives
\begin{equation}\label{eq:overlap}
 S<4n/3,\qquad \Delta_{\ge T}<n/\log T.
\end{equation}
The same charging, with charge $2s$ instead, bounds all deleted edges
in a root of order $m_0$ by $10sm_0\log m_0$.

For a leaf $Y$, let $Y^\circ$ be its vertices occurring in no other leaf.
Fewer than $\tau$ deleted edges join any given vertex $h$ to $Y^\circ$.
If $h\in Y$, none do. Otherwise all these edges are cut at the last
node on the path to $Y$ containing $h$. There $h$ is on the opposite
exclusive side from $Y^\circ$, and the rule moving high-degree vertices
gives the bound. This also holds for a composite tree whose initial
splits have $s=0$, since those delete no edges.

\emph{A short long-cycle argument.}
An $(\eps,0)$-robust graph of sufficiently large order $N$ has a cycle
of length at least $\eps N/(8\log^2N)$. It is connected, since a
smallest component would otherwise contradict expansion.
Run depth-first search and put $k=\lfloor N/2\rfloor$. Choose a deepest
vertex $v$ whose subtree has more than $k$ vertices. Its child subtrees
have size at most $k$ and total size at least $k$. Choose a union $U$
of these subtrees with $k/2\le|U|\le k$: take one large enough child,
or combine smaller children greedily. Undirected depth-first search
has no cross edges, so every external neighbour of $U$ lies on the
root-to-$v$ path. There are $q\ge\eps k/(2\log^2N)$ such neighbours.
If $x$ is the highest one and $xu$ an edge to $U$, the tree path from
$x$ to $u$, followed by $ux$, is a cycle of length at least $q+1$.
For sufficiently large $N$, $q\ge2$, so this is a genuine cycle.

\emph{One round.}
First remove cycles of length at least $d\lambda^2$ until none remains.
Split the rest with $s=0$ into weak expanders. The preceding argument
bounds every resulting piece by $M$: otherwise, for large $D$,
\[
 c\frac{d^2}{\log^2(d^2)}>d\log^2d
\]
contradicts the absence of long cycles (use monotonicity of
$x/\log^2x$ for large $x$). Split each weak piece of order at least $P$
again with the $s,\tau$ in \eqref{eq:scales}, stopping also at order $<P$.
The condition for \eqref{eq:splitshape} holds since
$256s\log^2M\le4608\lambda^{102}<\tau$.
The weak and robust stages form one composite splitting tree, so its
total leaf order is still less than $4n/3$.
Every large terminal leaf $Y$ has a robust spanning graph $H_Y$;
we retain the full vertex set of every such leaf.

First assign all edges of $H_Y$ to their own leaf. Then assign every
still unassigned edge having both ends in any retained leaf to one such
leaf, using a fixed rule. Denote the assigned set by $E_r(Y)$.
Pass all remaining edges to the next round. This assignment implies
\begin{equation}\label{eq:separation}
 H_Y\subseteq E_r(Y),\qquad
 \text{no later edge has both ends in }Y.
\end{equation}
All assigned sets, and consequently all reservoirs $H_Y$, are
edge-disjoint over the entire hierarchy.

Only cut edges and edges in leaves of order $<P$ can remain. By
\eqref{eq:overlap} and the deleted-edge bound their average degree satisfies
\begin{equation}\label{eq:recursion}
 d_{r+1}\le160\lambda_r^{101}+\tfrac43P_r\le\lambda_r^{105}.
\end{equation}
Stop below $D$. The final residual costs at most $Dn/2$ single edges.
Round $r$ removes at most $n/(2\lambda_r^2)$ long cycles.

\emph{The only scale estimates needed.}
After increasing $D$, the sequences $M_r$, $\lambda_r$ and
$\mu_r:=\log\lambda_r$ each decrease by a factor at least two.
For $l\ge3$, putting $x=\lambda_{l-2}$ gives
\begin{equation}\label{eq:tworound}
 M_l\le2(105\log x)^{210}.
\end{equation}
Thus $M_l^a=o(x^b)$ for every fixed $a,b>0$, uniformly in $l$.
If $R$ is the last active round, the elementary estimate
\begin{equation}\label{eq:sums}
 \sum_r\frac{(R-r+1)^k}{\mu_r^a}
 \le\frac{\sum_{j\ge0}(j+1)^k2^{-aj}}{(\log\log D)^a}
\end{equation}
holds for fixed $a>0,k\ge0$. The analogous estimate holds with $\lambda$;
also $\sum_r M_r^{-a}\le2D^{-2a}$ for $a\ge1$.
These replace all iterated-logarithm and tower bookkeeping.

\emph{Roles and concentration.}
Let $A_r$ be the vertices lying in at least two retained leaves in round
$r$. Call them hubs. All other vertices of a retained leaf are unique
ports. A unique port in round $l$ is fresh if it lies in no retained
leaf from a round $r\le l-2$. Otherwise fix one such leaf $Y(u)$ as
its class, before any randomness. A vertex can be fresh only at its
first appearance and in the following round, so
\begin{equation}\label{eq:fresh}
 \sum_l|\mathrm{Fresh}_l|\le2n.
\end{equation}
Let $\mathrm{Dup}_r$ denote vertices duplicated anywhere in the full
splitting tree. A duplicated vertex in a retained leaf has been
duplicated at a node of order at least $P_r$. If it occurs in $c$
retained leaves, then $c\le2k$, where $k\ge1$ is its number of such
duplications. Consequently
\begin{equation}\label{eq:hub}
 \sum_Y|Y\cap\mathrm{Dup}_r|\le2n/\log P_r,
 \qquad \sum_{r,Y}|Y\cap A_r|=o_D(n).
\end{equation}
Here and below $o_D(n)$ means $\eta(D)n$ for an absolute function
$\eta(D)\to0$, uniformly in the graph and its number of rounds.

Let $m_{Y,l}$ be the maximum, over vertices $h$, of the number of
round-$l$ assigned edges $hu$ with $Y(u)=Y$, where $Y$ is from round $r$.
By \eqref{eq:separation}, $h\notin Y$. Ports outside $\mathrm{Dup}_r$
give cut edges of the round-$r$ tree and hence contribute fewer than
$\tau_r$ at $h$. The others contribute at most their number.
There are at most $2n/P_r$ retained leaves at round $r$, so
\begin{align}\label{eq:concentration}
 \sum_{Y,l}m_{Y,l}
 &\le2n\sum_r(R-r)\left(\frac{\tau_r}{P_r}
                         +\frac1{\log P_r}\right)=o_D(n).
\end{align}
We used $\tau_r/P_r=O(1/\lambda_r)$ and \eqref{eq:sums}.
Finally, if $\nu_l$ is the number of available classes at round $l$,
\begin{equation}\label{eq:ancestors}
 \nu_l\le4n/P_{l-2},\qquad
 \sum_{l\ge3}\nu_l M_l^a=o_D(n)\quad\text{for every fixed }a>0.
\end{equation}

\section{Retiring an arbitrary set of vertices}
\begin{lemma}[Retirement]\label{lem:retire}
There are absolute constants $C,N_0$ such that the following holds.
Let $O$ be an $(\eps,s)$-robust spanning graph on a set $Z$ of order
$N\ge N_0$, with $1/32\le\eps\le1$ and $s\ge C\log^{42}N$.
For every graph $E$ containing $O$ on $Z$, and every $P\subseteq Z$,
one can cover all edges of $E[P]$, together with some other edges of $O$,
by at most $200|P|$ edge-disjoint cycles and single edges.
Every remaining edge of $E$ has an endpoint outside $P$.
\end{lemma}
\begin{proof}
The case $P=\varnothing$ is immediate. Put $L=\log N$ and $m=\lceil L^6\rceil$.
First choose $m$ distinct neighbours $A(w)$ of each vertex $w$ in $O$,
so that every vertex belongs to at most
$b=\lceil2mL^2/\eps\rceil=O(L^8)$ of these sets.
Here is a Hall proof that respects distinctness. On the left put each
oriented edge $(w,u)$. On the right put $b$ load slots at $u$ and
$\deg_O(w)-m$ spare slots at $w$, joining $(w,u)$ to these slots.
For any collection of oriented edges, let $R$ be its second coordinates
and $S'$ the first coordinates whose demands exceed their spare slots.
Each excess is at most $m$, and such a $w$ has at most $m-1$ neighbours
outside $R$. If $S'=\varnothing$, Hall is immediate. Otherwise choose
$U\subseteq S'$ with $|S'|/2\le|U|\le2N/3$.
Delete its edges to outside $R$, using at most $(m-1)|U|$ deletions.
Expansion gives $|R|\ge\eps|S'|/(2L^2)$, so $b|R|\ge m|S'|$ pays
all excesses. Hall now matches all oriented edges. At least $m$ edges
out of each $w$ use load slots; choose any $m$ of their distinct heads.

Let $J=\lfloor\log(L/6)\rfloor$, so $6/L\le2^{-J}<12/L$.
Split $O$ randomly into $2J+1$ graphs: a reserve $M$ and connectors
$R_{j,c}$, $0\le j<J$, $c=1,2$.
Independently, let each $p\in P$ survive to level $j$ with probability
$2^{-j}$, in nested levels; vertices outside $P$ survive always.
Give every vertex an independent label $0,1,2$ with probabilities
$1/2,1/4,1/4$. Let $U_j$ be the survivors,
$P_j=P\cap U_j$, $W_j=P_j\setminus P_{j+1}$, and
\[
 Z_j=U_{j+1}\cap\{\text{label }0\},\qquad
 V_{j,c}=U_{j+1}\cap\{\text{label }c\}.
\]
We can fix an outcome with the following properties: each $R_{j,c}$
links every pair multiset of multiplicity at most $b+2$ through
$V_{j,c}$; each $w\in P$ has four reserve edges into $A(w)\cap Z_j$
for every $j<J$; and
\begin{equation}\label{eq:retire-sum}
 \sum_{j<J}|P_j|\le8|P|,\qquad |P_J|\le48|P|/L.
\end{equation}
To see existence, edge splitting gives robustness at least
$s/(2(2J+1))$ in every connector. Each vertex belongs independently
to $V_{j,c}$ with probability at least $1/L$; apply
Lemma~\ref{lem:link} by thinning, with $t=b+2=O(L^8)$.
The lemma requires only $O(L^{29})$ robustness; the assumed $CL^{42}$
is ample. All connector failures have total probability $o(1)$.
For each reserve requirement the expected number of candidates is at
least $3L^4$, with independent indicators, so Chernoff and a union bound
over $NJ$ requirements also give failure $o(1)$.
Each inequality in \eqref{eq:retire-sum} fails with probability at most
$1/4$, by Markov. Thus all events hold simultaneously for large $N$.

Process $W_0,W_1,\ldots,W_{J-1}$. At step $j$, the unused edges of
$E[P]$ are supported on $P_j$. Let $F_j$ be those meeting $W_j$.
At each $w\in W_j$ reserve four edges into $A(w)\cap Z_j$, two for
each phase $c=1,2$, and remove these from $F_j$ if present.
Give each remaining edge of $F_j$ a phase whose $V_{j,c}$ avoids both
its ends. This is possible: one end is in $W_j$, and there is at most
one other end in $U_{j+1}$.

Decompose each phase into at most $|P_j|$ paths, with each vertex an
endpoint at most twice, using the published path corollary.
At an endpoint in $W_j$ append a distinct phase-reserved edge.
Output unused reserve edges singly. Each extended path avoids
$V_{j,c}$ and has both endpoints in $U_{j+1}$.
Only its two appended endpoints can repeat interior vertices. Excising
these repetitions gives at most two simple cycles and one simple path
with the original endpoints. If these endpoints coincide, only cycles
remain. Any residual path has length at least two: every edge meets
$W_j$, whereas its endpoints avoid $W_j$.

The residual paths have endpoint multiplicity at most $b+2$.
There are at most two inherited endpoints at a vertex; all new ones
come from distinct reserved edges $wx$ with $x\in A(w)$, and there
are at most $b$ such $w$. Connect each residual path's ends through
$V_{j,c}$. It avoids its connector's interior, so their union is a
simple cycle. Connectors within a phase are edge-disjoint; colours
separate the two phases.

No edge is spent twice. Earlier steps use only edges meeting their own
$W_i$ or belonging to their own connector colour. They cannot use a
reserve edge from $W_j$ into $U_{j+1}$, or a current connector edge with
both ends in $U_{j+1}$. Connectors belonging to $E[P]$ are immediately
removed from its unused graph. The induction on support therefore
continues with $P_{j+1}$.

The steps cost at most $6\sum_j|P_j|+4\sum_j|W_j|\le52|P|$ objects.
The final graph is empty if $h\le1$. Otherwise it is on
$h\le48|P|/L$ vertices and has a decomposition
into at most $2h\log h+h\le144|P|$ objects.
For completeness, this elementary bound follows by removing cycles
greedily. A graph with $m>h$ edges has a subgraph of minimum degree
at least $m/h$, by iterative deletion of vertices of smaller degree.
A longest path in that subgraph gives a cycle of length at least $m/h$.
In each range $(2^ih,2^{i+1}h]$, at most $2h$ removals suffice; finally
pay the at most $h$ remaining edges. Thus the total is below $200|P|$.
All additional edges used belong to $O$, as required.
\end{proof}

\section{Closing two-edge paths into cycles}
\begin{lemma}[Greedy grouping]\label{lem:cherry}
Let $B$ be an edge set such that every edge has a port endpoint. Partition
the ports into classes $P_Y\subseteq Y$, and suppose every $B$-edge
at a port in $P_Y$ has its other endpoint outside $Y$.
Let $d,K$ be positive integers, with every port of degree at most $d$. Put
$m_Y=\max_x|\{u\in P_Y:xu\in B\}|$.
For each $Y$ suppose $K$ edge-disjoint available graphs can link every
pair multiset in $P_Y$ of multiplicity at most $d$, through pairwise
disjoint sets $T_{Y,1},\ldots,T_{Y,K}\subseteq Y\setminus P_Y$.
Reservoir edges are disjoint across all classes and avoid $B$.

Then $B$, together with some reservoir edges, decomposes into a residual
set $J\subseteq B$ and at most
\begin{equation}\label{eq:cherry-cost}
 \tfrac12\sum_Ym_Y+2d\,\#\{Y\}+|B|/(2K)
\end{equation}
simple cycles. For every nonport $x$ and class $Y$, $J$ has at most
one edge joining $x$ to $P_Y$.
\end{lemma}
\begin{proof}
Assign each edge to one port endpoint, arbitrarily. Within each class,
delete one edge at every centre of odd assigned degree, putting it in
$J$. Pair all other edges at each centre into two-edge paths $u-x-v$.
One such path conflicts with at most $m_Y/2+2d-3$ others by sharing a
vertex. Greedily put the paths into groups of size at most $K$, opening
a new group only if every old one is full or has a conflict. For a
class with $N_Y>0$ paths, this uses at most
$m_Y/2+2d-3+\lceil N_Y/K\rceil$ groups.
Indeed at most the conflict degree old groups can be forbidden by
conflicts, and every other old group must be full.

Order a group's vertex-disjoint paths as
$p_1-x_1-q_1,\ldots,p_k-x_k-q_k$ and connect $q_j$ to $p_{j+1}$
cyclically through $T_{Y,j}$. For each slot $j$ make all requests
simultaneously over all groups. A port occurs in at most $d$ requests,
so the required connectors exist. Within a group their interiors are
disjoint and avoid all the two-edge paths. Their cyclic concatenation
is therefore a simple cycle, even for a one-path group.
All cycles are edge-disjoint by the routing and reservoir hypotheses.
Summing the grouping bound proves \eqref{eq:cherry-cost}.
Every edge at a nonport was necessarily assigned toward its port, so
only the one parity deletion per centre and class can survive in $J$.
\end{proof}

\section{Moving the leftovers into a small graph}
We record the quotient argument with its exact lifting obligation.
Its purpose is to count copies of a vertex without ever identifying two
occurrences of that vertex inside one lifted cycle.

\begin{lemma}[Compression]\label{lem:quotient}
Suppose the vertices involved have disjoint roles: hubs $A$, ports $U$,
middles $X$, and pools $W_L,W_R$. There are edge-disjoint cores of two
kinds: hub edges $au$, and port objects $uv$ or $uxv$ with $u\ne v$.
At a port, the total number of hub edges and port-object ends is at
most an integer $d\ge2$; each middle lies in at most $d$ objects.
Every port $u$ has available donor neighbours
$C_L(u)\subseteq W_L$, $C_R(u)\subseteq W_R$, each of size at least
$H\ge4d+4$. Donor edges avoid all cores.
Suppose, for each hub $a$ and pool vertex $w$,
\begin{equation}\label{eq:mu-condition}
 |\{u:au\text{ is a core},\ w\in C_L(u)\cup C_R(u)\}|\le\mu(w).
\end{equation}
If the numbers of hub cores and port objects are each at most $nd$,
there are paid core edges and a simple graph $Q$ such that
\begin{equation}\label{eq:quotient-size}
 \#\mathrm{paid}+|V(Q)|
 \le |A|+d^3\sum_{w\in W}\mu(w)
       +3dt_0|W|+C n\left(\frac dH+\frac1d\right),
 \quad t_0=\lceil2\log d\rceil,
\end{equation}
where $W=W_L\sqcup W_R$. Every cycle/edge decomposition of $Q$ lifts
to a decomposition of all unpaid cores plus some donor edges using
at most twice as many objects. Donor edges assigned to singleton
quotient edges are left unused.
\end{lemma}
\begin{proof}
For a loopless multigraph with edge multiplicities $m(v,w)$, put the
$i$th edge of each parallel class in a separate simple graph and delete
isolated vertices. The resulting disjoint union has
\begin{equation}\label{eq:ranks}
 \sum_v\max_w m(v,w)
\end{equation}
vertices. Each layer contains at most one copy of each real vertex.

\emph{Hub cores.}
At each hub group its edges in groups of size $\lfloor H/2\rfloor$,
except possibly the last, which can be smaller.
If there are $E$ hub edges, there are at most $|A|+4E/H$ groups.
Colour each group independently and uniformly with one of $d^3$ colours.
Pay an edge if another edge at its port has the same colour.
Edges at a port come from distinct hubs, hence distinct independently
coloured groups. The expected payment is at most $E(d-1)/d^3$.

For every surviving $au$, choose a donor neighbour $w$ greedily, avoiding
earlier choices at $u$ and earlier choices in its group.
Fewer than $d+H/2$ vertices are forbidden, so this is possible.
In each colour, replace $au$ by the quotient edge $aw$, representing
the path $a-u-w$, and split parallel edges into ranks.
For any hub, its multiplicity at a fixed pool vertex is at most its
number of groups in that colour. Thus total hub copies are at most
$|A|+4E/H$. By \eqref{eq:mu-condition}, total pool copies here are at
most $d^3\sum_W\mu(w)$.
Within a colour each port is used once, so internal vertices of the
representing paths are private.

\emph{Port objects.}
Greedily colour these objects with at most $3d$ colours so that objects
in one colour have disjoint ports and middles. This works because an
object conflicts with at most $3(d-1)$ others.
Orient every object between its two ports. At each port independently,
inject its tail ends uniformly into $C_L(u)$ minus the hub choices,
and its head ends uniformly into the corresponding right set.
Each target set has size at least $H-d\ge H/2$, and each domain size
at most $d$. All donor edges at a port are now distinct.
A port object becomes a quotient edge from its left to its right
junction. The representing path has length three or four.
For each colour, rank-split this bipartite multigraph as above.

Here is the entire load estimate. For independent labels
$Z_1,\ldots,Z_m$ with every atom of probability at most $p$, let $L_z$
be their counts. For each integer $t\ge1$,
\begin{equation}\label{eq:occupancy}
 \E\max_zL_z\le t\mathbf1_{m>0}+2emp+4e2^{-t}m.
\end{equation}
Indeed, put $\mu_z=\E L_z$. Factorial moments give
$\Prob(L_z\ge j)\le(e\mu_z/j)^j$. Set $T=\max(t,\lceil2emp\rceil)$.
For $j\ge T$, this is at most
$(e\mu_z/j)2^{1-j}$. Summing the tail gives
$\E[L_z\mathbf1_{L_z\ge T}]\le4e\mu_z2^{-T}$.
Now use $\max_zL_z\le T-1+\sum_z L_z\mathbf1_{L_z\ge T}$ and
$\sum_z\mu_z=m$. The case $m=0$ is immediate.

Fix an object colour and expose its tail junctions. For each left
junction $w$, the head junctions of objects starting there are
independent, with atom probabilities at most $2/H$: within this colour
all ports are distinct. Apply \eqref{eq:occupancy}, and then sum over
$w$ and colours. Reverse left and right for the other side.
If there are $P$ port objects, \eqref{eq:ranks} gives expected copies at most
\[
 3dt_0|W|+8eP/H+8e2^{-t_0}P.
\]
This bound holds conditional on every outcome of the hub step.
The expectation of the \emph{single sum} of payments and all copies
therefore satisfies \eqref{eq:quotient-size}. Averaging fixes one outcome
with that bound.

\emph{Lifting.}
Keep different colours, ranks and kinds in separate components of $Q$.
Within each component, distinct quotient vertices represent distinct real
branch vertices. Representing paths have private interior vertices, and
no interior vertex is a branch vertex.
Thus a simple cycle of $Q$ lifts to a simple cycle. All representing
paths are edge-disjoint, even across components, because cores and
donor choices are distinct. For a singleton quotient edge, output only
its core: one edge, or two edges for $uxv$. Return its donor edges
unused. This proves the precise lifting assertion and the factor two.
\end{proof}

\section{Preparing the reservoirs}
All hierarchy choices and port classes are now fixed.
For a leaf $Y$ from round $r$ put
\[
 k_r=\sum_{l\ge r+2}(M_l^2+1).
\]
By \eqref{eq:tworound}, $k_r=\lambda_r^{o(1)}$. If $k_r>0$, colour
each edge of $H_Y$ independently: an own colour with probability
$1/2$, and each of $k_r$ lending colours with probability
$p_r=1/(2k_r)\ge\lambda_r^{-4}$.
For each target round $l$, reserve $M_l^2$ colours for closing paths
and one colour for quotient donor edges. If $k_r=0$, keep all edges
in the own colour.

For each pair $(Y,l)$, independently give every vertex of $Y$ one of
$M_l^2$ slot labels, each with probability $\rho_l=M_l^{-4}$, or a
neutral label with the remaining probability. Write $T_{Y,l,j}$ for
the slot sets. They are disjoint for fixed $(Y,l)$, and are independent
of edge colours. Call $Y$ good if its own graph is robust enough for
Lemma~\ref{lem:retire}, all lending graphs have the robustness from
Lemma~\ref{lem:split}, and each closing colour links demands of
multiplicity $M_l$ through its corresponding $T_{Y,l,j}$.

These conditions fail with probability at most $|Y|^{-2}$.
Here are the scale checks, to make this assertion quantitative.
Writing $N=|Y|$, we have
\[
 \lambda_r^{103}\le N\le M_r,\qquad \log N\le3\lambda_r,
 \qquad M_l\le\lambda_r\quad(l\ge r+2).
\]
An atom has robustness at least $\lambda_r^{96}/2$ by direct Bernoulli
edge splitting. The linking requirement is only
$C M_l^{13}(\log N)^{18}\le C'\lambda_r^{31}$.
The own graph has robustness at least $s_r/4\gg C\log^{42}N$.
There are at most $k_r+1\le\lambda_r^4+1$ graphs, so their splitting
failures and $N^{-3}$ linking failures sum to at most $N^{-2}$ for
large $D$. All labels and class choices were fixed independently in
the required order.

A classed port $u$ at round $l$ is \emph{lost} if its class $Y(u)$ is
bad, or if it has a nonneutral slot label for $(Y(u),l)$.
Write $\Lost_l$ for these ports. Since the class was fixed before sampling,
\begin{equation}\label{eq:lost}
 \Prob(u\in\Lost_l)\le M_l^{-2}+|Y(u)|^{-2}\le2M_l^{-2},
 \qquad \E\sum_l(200+M_l)|\Lost_l|=o_D(n).
\end{equation}
A surviving port avoids every slot set of its class. Bad leaves lend
nothing: every possible consumer of such a class has been declared lost.

\emph{Disjoint pools for the quotient.}
All pool quantities and their sums below are indexed by $3\le l\le R$;
there are no lost ports in the first two rounds.
Independently of the preceding experiment, give each original vertex
one label $(l,r,\sigma)$, where $3\le l\le R$, $r\le l-2$ and
$\sigma\in\{L,R\}$, with probability
\[
 \tfrac12q_l\pi_{l,r},\qquad q_l=M_l^{-5},\quad
 \pi_{l,r}=2^{-(l-1-r)}.
\]
The total mass is less than one by \eqref{eq:sums}; all remaining mass
gets a null label. Let $W_{l,r,\sigma}$ be these disjoint sets,
$W_{l,\sigma}=\bigcup_rW_{l,r,\sigma}$, and
$W_l=W_{l,L}\sqcup W_{l,R}$.
For a possible classed port $u$ with class $Y$ from round $r$, its
candidates $C_\sigma(u)$ are its neighbours in $W_{l,r,\sigma}$ by
edges of $Y$'s quotient lending colour for round $l$.

For each side, the candidate count is a sum of independent Bernoulli
variables over distinct neighbours. Robustness gives
$\deg_{H_Y}(u)>s_r$, so its mean is at least
\[
 \lambda_r^{96}2^{-(l-r)}M_l^{-5}
 \ge\lambda_{l-2}^{95}/(4M_l^5).
\]
For the last inequality, set $j=l-2-r$ and use
$\lambda_r\ge2^j\lambda_{l-2}$.
Set
\begin{equation}\label{eq:H}
 H_l=\left\lfloor\frac{\lambda_{l-2}^{95}}{16M_l^5}\right\rfloor.
\end{equation}
By \eqref{eq:tworound}, $H_l\ge M_l^{10}$ for large $D$.
Call a possible port deficient if either candidate set has size less
than $H_l$. Chernoff bounds this probability by $2e^{-H_l/4}$.

\emph{Accounting that allows the later graph to depend on the pools.}
For a hub $h\in A_l$, let $c_l(h)$ be the number of classes with any
possible assigned round-$l$ edge from $h$ to a port of that class.
Define a deterministic weight
\[
 \omega_l(v)=\begin{cases}c_l(v)&v\in A_l,\\M_l&v\notin A_l.
 \end{cases}
\]
After Lemma~\ref{lem:cherry}, any residual $J_l$ has hub degrees at
most $c_l(h)$ and nonhub degrees at most $M_l$. Thus the number of
its edges touching $W_l$ is always at most
$\sum_{v\in W_l}\omega_l(v)$, even when $J_l$ depends on the pools.
Each hub--class incidence has a distinct original edge witnessing it,
and there are at most $nM_l$ edges touching possible ports. Hence
\begin{equation}\label{eq:pool-payment}
 \sum_v\omega_l(v)\le2nM_l,\qquad
 \E\sum_{v\in W_l}\omega_l(v)\le2n/M_l^4.
\end{equation}
Similarly all residual edges at deficient ports are dominated by $M_l$
times the number of deficient possible ports; its expectation is at
most $2nM_le^{-H_l/4}=o(n/M_l^4)$.

Let $\mu_r(w)$ count the round-$r$ retained leaves containing $w$.
By \eqref{eq:overlap}, $\sum_w\mu_r(w)<4n/3$.
For $w\in W_{l,r,L}\cup W_{l,r,R}$ use $\mu(w)=\mu_r(w)$.
The hub--class uniqueness of $J_l$ proves \eqref{eq:mu-condition}:
hub edges that can choose $w$ have distinct classes from round $r$,
all containing $w$.
Put $t_l=\lceil2\log M_l\rceil$ and
\[
 T_l=M_l^3\sum_{w\in W_l}\mu_{r(w)}(w)+3M_lt_l|W_l|.
\]
Its expectation is at most
$\tfrac43n/M_l^2+3nt_l/M_l^4$.
Consequently the expectation of the single nonnegative sum
\begin{equation}\label{eq:objective}
 Z=\sum_{l=3}^R\left[(200+M_l)|\Lost_l|
       +\sum_{v\in W_l}\omega_l(v)
       +M_l\#\{\text{deficient possible ports in round }l\}+T_l\right]
\end{equation}
is $o_D(n)$. Fix one outcome with $Z\le\E Z$.
This one averaging choice controls all these terms together; no
independence from the adaptive residual graphs is asserted.

\section{Assembly and the contraction}
Process rounds in descending order. For a retained leaf $Y$, remove
from $E_r(Y)$ exactly the reservoir edges already borrowed and used by later rounds,
and call the remainder $F_Y$. If $Y$ is good, its own graph
$\Own_Y$ is wholly available in $F_Y$. If $Y$ is bad, it has lent
nothing, so its full $H_Y$ is available. Use this available robust
graph as the device for Lemma~\ref{lem:retire}.

Within $Y$ retire its hubs, its fresh unique ports and its lost ports.
The remaining edges $B_Y\subseteq F_Y$ all touch a surviving classed
port. The retirement cost summed over all pieces is
\[
 200\sum_{r,Y}|Y\cap A_r|+200\sum_r|\mathrm{Fresh}_r|
      +200\sum_r|\Lost_r|
 \le400n+o_D(n).
\]
For the first two rounds all unique ports are fresh, so retirement
already finishes those rounds.

At round $l\ge3$ apply Lemma~\ref{lem:cherry} to
$B=\bigcup_Y B_Y$ with $d=M_l$ and $K=M_l^2$.
The separation property \eqref{eq:separation} places the centre of a
class-$Y$ edge outside $Y$. Surviving class-$Y$ ports avoid all of
$T_{Y,l,j}$. Good reservoirs therefore meet all hypotheses of the
lemma. The colours for this target round are still unused.
Since $|B|\le nM_l$, the cycle cost is at most
\[
 \tfrac12\sum_Ym_{Y,l}+2M_l\nu_l+n/(2M_l).
\]
Its sum is $o_D(n)$ by \eqref{eq:concentration},
\eqref{eq:ancestors} and \eqref{eq:sums}.
Record every reservoir edge used by these closing cycles against its
unique owner, for removal when that owner is processed.

The resulting $J_l$ has every edge incident to a surviving classed
port, degree at most $M_l$ outside the hubs, and at most one edge at
each hub and class. Pay all its edges touching lost vertices, pool
vertices or deficient ports. These payments are bounded by the
corresponding terms of \eqref{eq:objective}. At every remaining fresh
vertex, pair its edges into two-edge port paths and pay an unpaired
edge, if any. The latter payments sum to at most $2n$ by
\eqref{eq:fresh}. The two ends of each fresh path are distinct because the graph
is simple. The cores now consist of hub edges, port--port edges and
port--fresh--port paths. Their roles are disjoint from $W_l$.

Apply Lemma~\ref{lem:quotient} with $d=M_l$, $H=H_l$. Both kinds
of core number at most $nM_l$. For every history, its averaging
argument produces paid cores and a simple graph $Q_l$ with
\begin{equation}\label{eq:Ql}
 \#\mathrm{paid}_l+|V(Q_l)|
 \le |A_l|+T_l+C n(M_l/H_l+1/M_l).
\end{equation}
Choose a minimum cycle/edge decomposition of $Q_l$ and lift it.
Record only donor edges actually used in lifted cycles. In particular,
assigned connectors belonging to singleton quotient edges are returned.
Remove the recorded edges from their owners when those owners are
processed later in the descending order.

All edges are accounted for exactly once. The assigned sets
$E_r(Y)$, the removed long cycles and the final sparse residual
partition the original edge set. Reservoir colours have exclusive
consumers. Each borrowed edge is deleted once from its owner's
assigned set; each unused donor stays with its owner. Local retirement,
closing, and quotient lifting are exact decompositions of their stated
edge sets. Thus the descending construction neither reuses nor loses
an edge. The estimates hold uniformly in the earlier choices of
quotient decompositions.

Summing \eqref{eq:Ql}, using \eqref{eq:objective}, \eqref{eq:hub},
$H_l\ge M_l^{10}$ and \eqref{eq:sums}, gives
\[
 \sum_{l=3}^R\bigl(\#\mathrm{paid}_l+|V(Q_l)|\bigr)=o_D(n).
\]
All other costs total at most $C_0(D)n$: retirement costs $400n+o_D(n)$,
unpaired fresh edges at most $2n$, long cycles and closing cycles
$o_D(n)$, pool/lost/deficient payments $o_D(n)$, and the final
residual $Dn/2$. Choose $D$ once so that $\sum_l|V(Q_l)|\le n/4$.
All requirements on $D$ involved only absolute constants and fixed
powers of logarithms, never $n$ or the desired final constant.
We have proved \eqref{eq:reduction}.

Finally induct on $n$. The empty graph is immediate. Every $Q_l$ has
fewer than $n$ vertices for $n>0$, so
\[
 f(G)\le C_0n+2\sum_l2C_0|V(Q_l)|\le2C_0n.
\]
This completes the reconstruction of Theorem~\ref{thm:main}.

\paragraph{Provenance and status.}
The hierarchy, lending strategy and reduction to small quotients come
from the submitted proof \cite{Claim}, building on \cite{BM}.
Keeping every piece whole, the bounded-degree sampling certificate,
the direct cycle grouping, and the simplified quotient are changes
to that argument. The main external combinatorial inputs are
Lov\'asz's path corollary and Haxell's theorem, stated in Section 2.
Independent static checks found no remaining gap, but the revised
proof has not been formally verified. No Lean compilation was run.

\begin{thebibliography}{9}
\bibitem{BM} M. Buci\'c and R. Montgomery,
\emph{Towards the Erd\H{o}s--Gallai cycle decomposition conjecture},
Advances in Mathematics 437 (2024), 109434.
\href{https://arxiv.org/abs/2211.07689}{arXiv:2211.07689}.
\bibitem{Haxell} P. E. Haxell,
\emph{A condition for matchability in hypergraphs},
Graphs and Combinatorics 11 (1995), 245--248.
\bibitem{Claim} R. Coffey, \emph{erdos-gallai-lean}, proof claim and
accompanying manuscript, inspected commit
\texttt{660ae11bf3e975fca77869d7f2cf227d05445140}.
\href{https://github.com/steelwheel01/erdos-gallai-lean}{Repository};
\href{https://www.erdosproblems.com/forum/thread/184/proof-claims\#proof-claim-379}{proof claim}.
\end{thebibliography}
\end{document}
